
\documentclass[11pt,oneside]{article}
\usepackage[a4paper,margin=27mm,headheight=15pt]{geometry}
\usepackage[T1]{fontenc}
\usepackage[utf8]{inputenc}
\IfFileExists{libertinus.sty}{\usepackage{libertinus}}{\usepackage{lmodern}}
\usepackage{microtype}
\usepackage{amsmath,amssymb,amsthm,mathtools,mathrsfs}
\usepackage{bm}
\usepackage{booktabs,longtable,array,tabularx}
\usepackage{graphicx}
\usepackage{pgfplots}
\usepackage{xcolor}
\usepackage{enumitem}
\usepackage{fancyhdr}
\usepackage{titlesec}
\usepackage{listings}
\usepackage{xurl}
\usepackage{hyperref}
\usepackage[nameinlink,noabbrev]{cleveref}
\definecolor{linkblue}{RGB}{25,75,135}
\definecolor{shadegray}{RGB}{246,247,249}
\definecolor{rulegray}{RGB}{195,201,208}
\pgfplotsset{compat=1.18}
\hypersetup{
  colorlinks=true,
  linkcolor=linkblue,
  citecolor=linkblue,
  urlcolor=linkblue,
  pdftitle={Fabius inverse and saddle research frontiers},
  pdfsubject={Unformalized inverse consequences and saddle-coefficient research obligations},
  pdfkeywords={Fabius function, inverse Fabius function, Rvachev up-function, Lambert W, saddle point, asymptotic expansion, Stirling numbers}
}
\pagestyle{fancy}
\fancyhf{}
\fancyhead[L]{\small Fabius inverse and saddle research frontiers}
\fancyhead[R]{\small\nouppercase{\leftmark}}
\fancyfoot[C]{\thepage}
\renewcommand{\headrulewidth}{0.4pt}
\titleformat{\section}{\Large\bfseries}{\thesection}{0.7em}{}
\titleformat{\subsection}{\large\bfseries}{\thesubsection}{0.7em}{}
\titleformat{\subsubsection}{\normalsize\bfseries}{\thesubsubsection}{0.7em}{}
\setcounter{secnumdepth}{3}
\setcounter{tocdepth}{2}
\setlist[itemize]{leftmargin=2em,itemsep=0.25em,topsep=0.4em}
\setlist[enumerate]{leftmargin=2.3em,itemsep=0.25em,topsep=0.4em}
\setlength{\emergencystretch}{2em}
\newtheorem{theorem}{Theorem}[section]
\newtheorem{proposition}[theorem]{Proposition}
\newtheorem{lemma}[theorem]{Lemma}
\newtheorem{corollary}[theorem]{Corollary}
\newtheorem{conjecture}[theorem]{Conjecture}
\theoremstyle{definition}
\newtheorem{definition}[theorem]{Definition}
\newtheorem{algorithm}[theorem]{Algorithm}
\newtheorem{example}[theorem]{Example}
\theoremstyle{remark}
\newtheorem{remark}[theorem]{Remark}
\newtheorem{warning}[theorem]{Warning}
\newcommand{\R}{\mathbb R}
\newcommand{\C}{\mathbb C}
\newcommand{\Q}{\mathbb Q}
\newcommand{\N}{\mathbb N}
\newcommand{\Z}{\mathbb Z}
\newcommand{\Up}{\operatorname{up}}
\newcommand{\sinc}{\operatorname{sinc}}
\newcommand{\wt}{\operatorname{wt}_2}
\newcommand{\e}{\mathrm e}
\newcommand{\ii}{\mathrm i}
\newcommand{\dd}{\,\mathrm d}
\newcommand{\Li}{\operatorname{Li}}
\newcommand{\Var}{\operatorname{Var}}
\newcommand{\E}{\mathbb E}
\newcommand{\Prob}{\mathbb P}
\newcommand{\bigO}{\mathcal O}
\newcommand{\smallo}{o}
\newcommand{\supp}{\operatorname{supp}}
\newcommand{\sgn}{\operatorname{sgn}}
\newcommand{\lcm}{\operatorname{lcm}}
\newcommand{\vtwo}{v_2}
\newcommand{\qbinom}[3]{\genfrac{[}{]}{0pt}{}{#1}{#2}_{#3}}
\newcommand{\repo}{\nolinkurl{Analysis/FabiusFunction}}
\newenvironment{proofidea}{\par\smallskip\noindent\textit{Proof architecture.}\ }{\hfill$\square$\par\smallskip}
\newenvironment{boxedremark}{%
  \par\smallskip\noindent\begin{center}\begin{minipage}{0.94\textwidth}%
  \color{black}\hrule\vspace{0.6em}\small
}{%
  \vspace{0.6em}\hrule\end{minipage}\end{center}\smallskip
}
\lstdefinestyle{pseudo}{
  basicstyle=\ttfamily\small,
  backgroundcolor=\color{shadegray},
  frame=single,
  rulecolor=\color{rulegray},
  columns=fullflexible,
  keepspaces=true,
  showstringspaces=false,
  xleftmargin=0.7em,
  xrightmargin=0.7em,
  aboveskip=0.8em,
  belowskip=0.8em
}

% Companion-specific notation, added after the retained house preamble.
\newcommand{\InvF}{G}
\newcommand{\clamp}{\operatorname{clamp}_{[0,1]}}
\newcommand{\Gauss}{\mathfrak G}
\newcommand{\Comp}{\mathcal C}
\newcommand{\Ham}{H}
\newcommand{\Formalized}{\textsc{Lean}}
\newcommand{\Derived}{\textsc{Derived}}
\newcommand{\Main}{\mathcal M}

\title{\bfseries Fabius Inverse and Saddle Research Frontiers\\[0.35em]
\large Derived consequences and coefficient formulas awaiting formalization}
\author{Research companion to the formal development\\
\small Vladimir Reshetnikov's \texttt{ProveIt} repository}
\date{Repository snapshot: 25 August 2026}

\begin{document}
\maketitle

\begin{abstract}
The Fabius function is unusually rich: a single compactly supported
$C^\infty$ density produces a self-similar distribution function, exact
dyadic arithmetic, Fourier products, a probability model, and a
small-argument expansion whose coefficients are smooth periodic functions.
The companion article \emph{The Fabius Function and Rvachev's Up-Function}
develops most of that landscape.  The present paper is deliberately
different.  It does not repeat the introductory theory.  Instead it fills
two identifiable gaps in the exposition by turning recently formalized
material into ordinary mathematics.

The first gap is the inverse.  We construct the totalized inverse
$\InvF=F^{-1}$ on the whole real line, prove its global monotonicity,
continuity, reflection law, comparison calculus, and exact values at all
outputs of the dyadic evaluator, and show how every flatness estimate for
$F$ becomes an effective steepness estimate for $\InvF$.  We then derive
the classical interior calculus of the inverse and invert the sharp
Lambert-phase asymptotic.  The result is the explicit endpoint law
\[
 \InvF(y)\sim
 \frac{\sqrt{\log(1/y)}}{\e\sqrt{\log 2}}\,
 \exp\!\left(-\sqrt{2\log 2\,\log(1/y)}\right)
 \qquad (y\downarrow0),
\]
together with its reflected analogue at $1$.  This explains in one formula
why the inverse is steeper than every positive power of $y$, although it
still tends to zero faster than every negative power of $\log(1/y)$.

The second gap concerns the internal machinery of the all-orders saddle
expansion.  We give complete derivations of the closed periodic jet
formula, its shifted signed-Stirling conversion, the collapse of the
centered exponent coefficients, the formal mass-to-logarithm recursions,
the leading-derivative law, and the full second logarithmic correction.
Finally, we give a proof-level account of the arbitrary-order analytic
remainder: the order-dependent central window, vertical Taylor control,
finite exponential substitution, Gaussian contraction, polynomial-tail
replacement, minor-arc estimates, logarithmic transfer, and return from
the dyadic logarithmic scale to $x\downarrow0$.

Statements explicitly represented by theorem declarations in the Lean
development are marked \Formalized.  Consequences proved here from those
statements by standard real or analytic arguments are marked \Derived.
\end{abstract}

\begin{boxedremark}
\textbf{Formalization boundary.}
This document is a research-frontier notebook, not part of the project's
formalization-backed primary exposition.  It deliberately preserves a mixture of proved
Lean inputs, ordinary mathematical deductions, symbolic coefficient calculations, and
numerical evidence so that the unformalized steps are not lost.  A statement here is a
project theorem only when an exact current Lean declaration is named for the full statement;
the labels \Formalized\ and \Derived\ record the draft's provenance but do not themselves
certify that every displayed consequence has been formalized.  In particular, H\"older
consequences beyond the now-formalized inverse-curvature identities, inverse endpoint
asymptotics, the identification with
standard shifted Stirling numbers, ordered-composition coefficient formulas, and the
logarithmic highest-derivative law remain research obligations.
\end{boxedremark}

\tableofcontents
\clearpage

\section{Purpose, scope, and status conventions}

\subsection{A supplement rather than a second survey}

Let $F$ denote the bounded, totalized Fabius function: $F(x)=0$ for
$x\leq0$, $F(x)=1$ for $x\geq1$, and on $[0,1]$ it is the unique strictly
increasing $C^\infty$ function associated with Rvachev's compactly
supported up-function.  We use only a small portion of the foundational
theory:
\begin{align}
 F(0)&=0,&F(1)&=1,&F(1-x)&=1-F(x), \label{eq:F-basic}\\
 F'(x)&=2\Up(2x-1),&&&0<x<1, \label{eq:Fprime-up}
\end{align}
where $\Up$ is positive on $(-1,1)$ and $\Up(0)=1$.  Consequently $F$ is a
continuous strictly increasing bijection of $[0,1]$ onto itself.

The main companion article already explains the random-series model,
infinite convolutions, the Fourier transform, moment recurrences,
Thue--Morse identities, dyadic rationality, and the broad shape of the
Lambert saddle analysis.  Repeating those chapters would obscure the
missing material.  The division of labor adopted here is summarized in
\cref{tab:audit}.

\begin{longtable}{>{\raggedright\arraybackslash}p{0.24\textwidth}
                  >{\raggedright\arraybackslash}p{0.35\textwidth}
                  >{\raggedright\arraybackslash}p{0.31\textwidth}}
\caption{Expository audit and the role of this companion.}
\label{tab:audit}\\
\toprule
Topic & Situation in the main article & Treatment here\\
\midrule
\endfirsthead
\toprule
Topic & Situation in the main article & Treatment here\\
\midrule
\endhead
Inverse function & The totalized inverse module and its endpoint estimates
are absent. & Full construction, exact identities, dyadic inversion,
effective steepness, endpoint limits, regularity, and inverse
asymptotics.\\
Closed saddle jets & Closed formulas are displayed, but the operator
calculus and the shifted Stirling conversion are compressed. & Complete
induction, forcing-term propagation, coefficient table, and the reason for
the index shift $s(n+1,m+1)$.\\
Mass and logarithmic coefficients & The all-orders result is stated and
some coefficients are listed. & Formal exponential recursion, Gaussian
contraction, composition formulas, structural bounds, and the leading
derivative law.\\
Second correction & The final polynomial appears, while much of the
order-four algebra is easy to miss. & A line-by-line calculation from
$E_1,\dots,E_4$ through $B_4$, Gaussian moments, $a_2$, and $A_2$.\\
Arbitrary-order remainder & The main article gives a high-level proof
architecture. & The order-dependent contour window and every analytic
error mechanism are separated and justified.\\
Probability, Fourier theory, dyadic arithmetic & Already developed in
substantial depth. & Used as input only; not duplicated.\\
\bottomrule
\end{longtable}

\subsection{What the labels mean}

A paragraph headed \Formalized\ restates a mathematical theorem whose
content is represented directly in the Lean files under
\begin{center}
\path{Analysis/FabiusFunction/Lean/FabiusFunction}.
\end{center}
A paragraph headed \Derived\ gives a paper-level consequence.  Its proof
uses formalized results plus a standard theorem such as the inverse
function theorem, ordinary asymptotic inversion, or elementary formal
power-series algebra.  The distinction matters especially in
\cref{sec:inverse-calculus,sec:coefficient-formulas}: those sections contain
new synthesis, but they do not pretend that every displayed corollary has
a declaration of its own.

\begin{boxedremark}
The exact Lambert phase is never silently replaced inside a periodic
coefficient.  If $\lambda=\lambda(x)$ is the exact lower solution of
$x=\lambda2^{-\lambda}$, then a term such as $A_j(\lambda)$ remains at that
exact phase.  Replacing $\lambda$ by an asymptotic approximation inside
$A_j$ is a separate composition problem, because a bounded phase error
does not become small modulo the period.
\end{boxedremark}

\section{The inverse as a global mathematical object}
\label{sec:inverse}

\subsection{The order isomorphism on the unit interval}

Set
\[
 I=[0,1],\qquad
 \clamp(y)=\min\{1,\max\{0,y\}\}.
\]
Because $F|_I$ is a continuous strictly increasing bijection $I\to I$, it
is an order isomorphism.  Its ordinary inverse will first be denoted by
$F_I^{-1}:I\to I$.

\begin{definition}[Totalized inverse]
\label{def:totalized-inverse}
The totalized inverse of the Fabius function is
\[
 \InvF(y)=F_I^{-1}\bigl(\clamp(y)\bigr),\qquad y\in\R.
\]
Thus the inverse is extended by constants outside its natural range,
exactly as $F$ itself is extended by constants outside its natural domain.
\end{definition}

\begin{theorem}[Global inverse package; \Formalized]
\label{thm:inverse-package}
The function $\InvF$ has the following properties.
\begin{enumerate}
\item $0\leq \InvF(y)\leq1$ for every $y\in\R$.
\item If $0\leq y\leq1$, then
\[
 F(\InvF(y))=y.
\]
If $0\leq x\leq1$, then
\[
 \InvF(F(x))=x.
\]
\item $\InvF$ is strictly increasing on $[0,1]$, monotone on $\R$, and
continuous on $\R$.
\item The restriction $\InvF|_I$ is a bijection $I\to I$.
\end{enumerate}
\end{theorem}

\begin{proof}
The clamp maps $\R$ continuously and monotonically into $I$.  The inverse
of an order isomorphism between compact real intervals is itself an order
isomorphism, hence is continuous and strictly increasing.  Their
composition is therefore continuous and monotone on the whole line and
takes values in $I$.  If $y\in I$, then $\clamp(y)=y$, so the usual
right-inverse identity for $F_I^{-1}$ gives $F(\InvF(y))=y$.  Likewise,
$F(x)\in I$ for $x\in I$, and the left-inverse identity gives
$\InvF(F(x))=x$.  Strict increase and bijectivity on $I$ are inherited
from $F_I^{-1}$.
\end{proof}

The totalization is not cosmetic.  It produces a globally continuous
monotone function and makes the reflection law valid without restricting
the variable to $[0,1]$.

\subsection{A four-way comparison calculus}

The order-isomorphism viewpoint turns inverse inequalities into direct
inequalities without any numerical inversion.

\begin{theorem}[Comparison equivalences; \Formalized]
\label{thm:inverse-comparisons}
For $x,y\in[0,1]$,
\begin{align}
 \InvF(y)\leq x&\iff y\leq F(x), \label{eq:inv-le}\\
 F(x)\leq y&\iff x\leq \InvF(y), \label{eq:F-le}\\
 \InvF(y)<x&\iff y<F(x), \label{eq:inv-lt}\\
 F(x)<y&\iff x<\InvF(y). \label{eq:F-lt}
\end{align}
\end{theorem}

\begin{proof}
Apply the strictly increasing function $F$ to the two sides of an
inequality involving $\InvF(y)$, then use $F(\InvF(y))=y$.  For example,
\[
 \InvF(y)\leq x
 \iff F(\InvF(y))\leq F(x)
 \iff y\leq F(x).
\]
The other three equivalences are identical, with the order reversed in
placement but not in direction because $F$ and $\InvF$ are increasing.
\end{proof}

\begin{remark}
The equivalences are often more useful than an explicit inverse formula.
They let one convert a known dyadic or asymptotic bound for $F$ into a
location bound for $\InvF$ with no loss in constants.
\end{remark}

\subsection{Tails, endpoints, and reflection}

\begin{theorem}[Constant tails and symmetry; \Formalized]
\label{thm:inverse-symmetry}
For every $y\in\R$,
\begin{align}
 y\leq0&\Longrightarrow \InvF(y)=0,\\
 y\geq1&\Longrightarrow \InvF(y)=1,\\
 \InvF(1-y)&=1-\InvF(y). \label{eq:inv-reflection}
\end{align}
In particular,
\[
 \InvF(0)=0,\qquad \InvF(1)=1,\qquad \InvF(1/2)=1/2.
\]
\end{theorem}

\begin{proof}
The two tail identities follow immediately from the clamp.  Suppose first
that $0\leq y\leq1$.  Then $\InvF(y)\in[0,1]$, and by the symmetry of $F$,
\[
 F(1-\InvF(y))
   =1-F(\InvF(y))
   =1-y.
\]
The point $1-\InvF(y)$ lies in $[0,1]$, so uniqueness of the inverse gives
$\InvF(1-y)=1-\InvF(y)$.  If $y\leq0$ or $y\geq1$, the same equation follows
from the two constant tails.  Setting $y=0,1,$ and $1/2$ gives the stated
special values.
\end{proof}

\section{Exact inversion of the dyadic evaluator}
\label{sec:dyadic-inverse}

\subsection{The exact statement}

Let $\mathsf F_n(a)$ denote the exact bounded dyadic evaluator from the
formal development: for $0\leq a\leq2^n$,
\[
 \mathsf F_n(a)=F\!\left(\frac{a}{2^n}\right),
\]
and for larger $a$ it is clamped at the endpoint value $1$.

\begin{theorem}[Dyadic inversion; \Formalized]
\label{thm:dyadic-inversion}
For all $n,a\in\N$,
\[
 \InvF\bigl(\mathsf F_n(a)\bigr)
 =\frac{\min\{a,2^n\}}{2^n}.
\]
In particular, if $0\leq a\leq2^n$, then
\[
 \InvF\!\left(F\!\left(\frac a{2^n}\right)\right)=\frac a{2^n}.
\]
\end{theorem}

\begin{proof}
When $a\leq2^n$, the argument $a/2^n$ lies in $I$, so the left-inverse
identity in \cref{thm:inverse-package} applies.  When $a\geq2^n$, the
bounded evaluator equals $1$, the minimum equals $2^n$, and both sides are
$1$.
\end{proof}

The theorem is stronger than a table of examples: every exact rational
produced by the dyadic algorithm immediately gives an exact algebraic
input for the inverse.

\subsection{A small exact table}

\begin{table}[ht]
\centering
\small
\begin{tabular}{@{}ccc@{}}
\toprule
$x$ & $F(x)$ & exact inverse statement\\
\midrule
$0$ & $0$ & $\InvF(0)=0$\\
$\dfrac1{16}$ & $\dfrac{143}{2073600}$
  & $\InvF\!\left(\dfrac{143}{2073600}\right)=\dfrac1{16}$\\[0.55em]
$\dfrac18$ & $\dfrac1{288}$
  & $\InvF\!\left(\dfrac1{288}\right)=\dfrac18$\\[0.55em]
$\dfrac14$ & $\dfrac5{72}$
  & $\InvF\!\left(\dfrac5{72}\right)=\dfrac14$\\[0.55em]
$\dfrac38$ & $\dfrac{73}{288}$
  & $\InvF\!\left(\dfrac{73}{288}\right)=\dfrac38$\\[0.55em]
$\dfrac12$ & $\dfrac12$
  & $\InvF\!\left(\dfrac12\right)=\dfrac12$\\[0.55em]
$\dfrac58$ & $\dfrac{215}{288}$
  & $\InvF\!\left(\dfrac{215}{288}\right)=\dfrac58$\\[0.55em]
$\dfrac34$ & $\dfrac{67}{72}$
  & $\InvF\!\left(\dfrac{67}{72}\right)=\dfrac34$\\[0.55em]
$\dfrac78$ & $\dfrac{287}{288}$
  & $\InvF\!\left(\dfrac{287}{288}\right)=\dfrac78$\\[0.55em]
$\dfrac{15}{16}$ & $\dfrac{2073457}{2073600}$
  & $\InvF\!\left(\dfrac{2073457}{2073600}\right)=\dfrac{15}{16}$\\
\bottomrule
\end{tabular}
\caption{Exact inverse values obtained from the dyadic evaluator and
reflection.}
\label{tab:inverse-dyadic}
\end{table}

\begin{remark}
The threshold $F(1/4)=5/72$ will reappear in the effective endpoint
steepness theorem.  Its occurrence is structural rather than arbitrary:
$1/4=2^{-2}$ is precisely the scale on which the quadratic flatness bound
is valid.
\end{remark}

\section{\texorpdfstring{Flatness of \(F\) becomes steepness of \(\InvF\)}{Flatness of F becomes steepness of its inverse}}
\label{sec:inverse-steepness}

\subsection{The transported power inequalities}

Write
\[
 C_n=2^{\binom{n+1}{2}}.
\]
The effective flatness theorem for $F$ says, in one of its standard forms,
that
\[
 0\leq x\leq2^{-n}
 \quad\Longrightarrow\quad
 F(x)\leq C_nx^n. \label{eq:flatness-input}
\]
The sharp version improves the right-hand side by the factor $n!$.

\begin{theorem}[Inverted effective and sharp flatness; \Formalized]
\label{thm:inverted-flatness}
Let $n\in\N$ and $0\leq y\leq1$.  If
\[
 2^n\InvF(y)\leq1,
\]
then
\begin{align}
 y&\leq C_n\InvF(y)^n, \label{eq:inverse-effective}\\
 y&\leq \frac{C_n}{n!}\InvF(y)^n. \label{eq:inverse-sharp}
\end{align}
Moreover, the scale condition follows from the argument condition
\[
 0\leq y\leq F(2^{-n}).
\]
\end{theorem}

\begin{proof}
Put $x=\InvF(y)$.  The scale hypothesis says $0\leq x\leq2^{-n}$, while
$F(x)=y$.  Substituting $x$ into the effective and sharp flatness bounds
for $F$ gives \eqref{eq:inverse-effective} and
\eqref{eq:inverse-sharp}.

Now suppose $y\leq F(2^{-n})$.  By
\eqref{eq:inv-le}, $\InvF(y)\leq2^{-n}$, which is exactly the required
scale condition.
\end{proof}

For $n\geq1$, taking positive $n$th roots yields the more geometric form.

\begin{corollary}[Root lower bounds; \Formalized\ input, \Derived\ form]
\label{cor:root-lower}
If $n\geq1$, $0<y\leq F(2^{-n})$, then
\begin{align}
 \InvF(y)&\geq C_n^{-1/n}y^{1/n}, \label{eq:root-effective}\\
 \InvF(y)&\geq
 \left(\frac{n!}{C_n}\right)^{1/n}y^{1/n}. \label{eq:root-sharp}
\end{align}
\end{corollary}

\begin{proof}
All quantities are nonnegative.  Rearrange
\cref{eq:inverse-effective,eq:inverse-sharp} and apply the monotonicity of
the positive $n$th-root function.
\end{proof}

\subsection{An effective bound against every line}

The quadratic case already detects infinite endpoint slope.

\begin{theorem}[Effective linear steepness; \Formalized]
\label{thm:effective-linear}
Let $M,y\in\R$.  If
\[
 0<y,\qquad y\leq F(1/4)=\frac5{72},\qquad
 8M^2y<1,
\]
then
\[
 My<\InvF(y).
\]
\end{theorem}

\begin{proof}
The comparison theorem gives $\InvF(y)\leq1/4$, so the quadratic inverse
flatness estimate applies:
\[
 y\leq8\InvF(y)^2. \label{eq:quadratic-inv}
\]
If $M\leq0$, then $My\leq0<\InvF(y)$ because $y>0$ and $\InvF$ is strictly
increasing on $[0,1]$.  Suppose $M>0$ and, toward a contradiction,
$\InvF(y)\leq My$.  Then
\[
 y\leq8\InvF(y)^2\leq8M^2y^2.
\]
Dividing by $y>0$ gives $1\leq8M^2y$, contrary to the last hypothesis.
\end{proof}

\subsection{Infinite one-sided slope}

\begin{theorem}[Endpoint quotient divergence; \Formalized]
\label{thm:inverse-quotient}
As $y\downarrow0$,
\[
 \frac{\InvF(y)}{y}\longrightarrow+\infty.
\]
By reflection, as $y\uparrow1$,
\[
 \frac{1-\InvF(y)}{1-y}\longrightarrow+\infty.
\]
\end{theorem}

\begin{proof}
Fix a real target $A$.  Put $B=|A|+1$.  For all sufficiently small
positive $y$ we have simultaneously
\[
 y<F(1/4),\qquad 8B^2y<1.
\]
By \cref{thm:effective-linear}, $By<\InvF(y)$, and hence
\[
 A\leq B<\frac{\InvF(y)}y.
\]
This is exactly the filter definition of divergence to $+\infty$.

For the right endpoint use \eqref{eq:inv-reflection}:
\[
 \frac{1-\InvF(y)}{1-y}
 =\frac{\InvF(1-y)}{1-y},
\]
and let $1-y\downarrow0$.
\end{proof}

\subsection{Failure of every positive Hölder estimate}

\begin{theorem}[Stronger than non-Lipschitz; \Derived]
\label{thm:no-holder}
For every $\alpha>0$,
\[
 \frac{\InvF(y)}{y^\alpha}\longrightarrow+\infty
 \qquad (y\downarrow0).
\]
Consequently there are no $C,\delta>0$ for which
\[
 \InvF(y)\leq Cy^\alpha\qquad(0\leq y\leq\delta).
\]
The reflected assertion holds at $1$:
\[
 \frac{1-\InvF(y)}{(1-y)^\alpha}\longrightarrow+\infty
 \qquad (y\uparrow1).
\]
\end{theorem}

\begin{proof}
Choose an integer $n>1/\alpha$.  On the window
$0<y\leq F(2^{-n})$, \cref{eq:root-effective} gives
\[
 \frac{\InvF(y)}{y^\alpha}
 \geq C_n^{-1/n}y^{1/n-\alpha}.
\]
The exponent $1/n-\alpha$ is negative, so the right-hand side tends to
$+\infty$.  A local Hölder upper bound would keep the quotient bounded and
is therefore impossible.  The statement at $1$ follows from reflection.
\end{proof}

\begin{boxedremark}
Each fixed $n$ supplies a root lower bound only on the shrinking window
$y\leq F(2^{-n})$.  There is no bound uniform in $n$.  This is not a
defect: it is exactly the scale dependence needed to prove that
$\InvF$ outruns \emph{every fixed} positive power.
\end{boxedremark}


\section{Classical calculus and the exact regularity of the inverse}
\label{sec:inverse-calculus}

This section uses the formalized inverse package together with standard
one-variable analysis.  Smoothness, the reciprocal first-derivative formula,
the exact second-derivative formula, and the strict closed-half shape are now
marked \Formalized; the general higher inverse-derivative formulas remain
\Derived.

\subsection{Smoothness and differential identities in the interior}

\begin{theorem}[Interior inverse calculus through first order; \Formalized]
\label{thm:inverse-calculus}
The inverse $\InvF$ is $C^\infty$ on $(0,1)$.  For $0<y<1$,
\begin{equation}
 \InvF'(y)
   =\frac1{F'(\InvF(y))}
     =\frac1{2\Up(2\InvF(y)-1)}>0.
 \label{eq:Gprime}
\end{equation}
\end{theorem}

\begin{proof}
For $x\in(0,1)$, \eqref{eq:Fprime-up} and positivity of $\Up$ on
$(-1,1)$ imply $F'(x)>0$.  The $C^\infty$ inverse function theorem
therefore applies at every interior point and gives a smooth local inverse.
Uniqueness identifies all these local inverses with $\InvF$.  Differentiating
$F(\InvF(y))=y$ and substituting \eqref{eq:Fprime-up} proves
\eqref{eq:Gprime}.
\end{proof}

The corresponding Lean declarations are \path{fabiusInv_contDiffOn_Ioo},
\path{fabiusInv_hasDerivAt},
\path{deriv_fabiusInv_eq_inv_two_mul_rvachevUp}, and
\path{deriv_fabiusInv_pos}.

\begin{proposition}[Second inverse derivative; \Formalized]
\label{prop:inverse-second-derivative}
For $0<y<1$,
\begin{equation}
 \InvF''(y)
   =-\frac{F''(\InvF(y))}
           {F'(\InvF(y))^3}.
 \label{eq:Gsecond}
\end{equation}
\end{proposition}

\begin{proof}
Differentiating the identity behind \eqref{eq:Gprime} once more gives
\[
 F''(\InvF(y))\InvF'(y)^2+
 F'(\InvF(y))\InvF''(y)=0.
\]
Substitution of \eqref{eq:Gprime} yields \eqref{eq:Gsecond}.  This is
\path{deriv_fabiusInv_hasDerivAt} and its pointwise corollary
\path{deriv_deriv_fabiusInv}.
\end{proof}

\begin{remark}[Higher inverse derivatives; \Derived]
Repeated differentiation expresses every higher derivative of $\InvF$ as a
rational polynomial in the derivatives of $F$ evaluated at $\InvF(y)$, with a
power of $F'(\InvF(y))$ in the denominator.  These general Fa\`a di Bruno
inverse-derivative polynomials are not separately packaged in the Lean API.
\end{remark}

\begin{corollary}[Midpoint differential data; \Derived]
\label{cor:midpoint-inverse}
At the fixed point $1/2$,
\[
 \InvF(1/2)=1/2,\qquad
 \InvF'(1/2)=\frac12,\qquad
 \InvF''(1/2)=0.
\]
\end{corollary}

\begin{proof}
The value follows from reflection.  Since $\Up(0)=1$,
\eqref{eq:Gprime} gives $\InvF'(1/2)=1/2$.  The up-function is even, so
$\Up'(0)=0$, hence $F''(1/2)=4\Up'(0)=0$; now use
\eqref{eq:Gsecond}.
\end{proof}

\subsection{Shape reversal under inversion}

The formal development proves that $F$ is strictly convex on
$(0,1/2)$ and strictly concave on $(1/2,1)$.  Inversion reverses these
two curvature signs.

\begin{proposition}[Concavity and convexity of the inverse; \Formalized]
\label{prop:inverse-shape}
The function $\InvF$ is strictly concave on $[0,1/2]$ and strictly convex
on $[1/2,1]$.
\end{proposition}

\begin{proof}
The inverse is strictly increasing and preserves the midpoint.  It therefore
maps the first half into itself, where $F'$ is strictly increasing and
positive.  Its reciprocal in \eqref{eq:Gprime} is strictly decreasing, so the
derivative criterion and continuity give strict concavity on the closed first
half.  On the second half $F'$ is strictly decreasing and positive, hence the
reciprocal derivative of the inverse is strictly increasing and the same
criterion gives strict convexity.

The closed-interval conclusions are
\path{strictConcaveOn_fabiusInv_firstHalf} and
\path{strictConvexOn_fabiusInv_secondHalf}.  Their derivative criteria inspect
only the two open interval interiors, while global continuity supplies the
endpoints; no finite endpoint derivative is asserted.
\end{proof}

\begin{remark}[Endpoint slope; \Derived]
The secant slopes from the origin diverge by \cref{thm:inverse-quotient}, and
reflection gives the right-endpoint counterpart.  Interpreting this together
with concavity as an extended one-sided derivative statement remains ordinary
analysis rather than a separately named Lean theorem.
\end{remark}

\subsection{\texorpdfstring{The \(C^\infty\) and analytic loci}{The smooth and analytic loci}}

\begin{theorem}[Exact regularity locus; \Derived]
\label{thm:inverse-regularity}
The totalized inverse has the following regularity.
\begin{enumerate}
\item $\InvF$ is $C^\infty$ precisely on
\[
 \R\setminus\{0,1\}.
\]
\item It has no finite derivative at $0$ or at $1$.
\item It is real analytic precisely on
\[
 \R\setminus[0,1].
\]
\end{enumerate}
\end{theorem}

\begin{proof}
On $(-\infty,0)$ and $(1,\infty)$ the function is constant, hence smooth
and analytic.  On $(0,1)$ smoothness follows from
\cref{thm:inverse-calculus}.  At $0$, the left derivative is $0$ because
the function is constant to the left, whereas the right difference
quotient $\InvF(y)/y$ tends to $+\infty$.  Thus no finite derivative
exists.  Reflection gives the same conclusion at $1$.

It remains to exclude analyticity in the open unit interval.  Suppose
$\InvF$ were analytic near some $y_0\in(0,1)$.  Its derivative there is
nonzero by \eqref{eq:Gprime}; the analytic inverse function theorem would
then make its local inverse $F$ analytic near $x_0=\InvF(y_0)$.  But the
Fabius function is nowhere analytic on the interior of its transition
interval.  This contradiction proves the claim.
\end{proof}

\section{Inverting the sharp small-argument asymptotic}
\label{sec:inverse-asymptotic}

The power inequalities show qualitative infinite steepness.  The sharp
Lambert expansion gives the actual scale.  This section extracts that
scale without needing the detailed value of the periodic fluctuation.

\subsection{The exact lower Lambert phase}

Put
\[
 L=\log2.
\]
For sufficiently small $x>0$, let $\lambda=\lambda(x)$ be the large
positive solution of
\begin{equation}
 x=\lambda2^{-\lambda}
   =\lambda\e^{-L\lambda}. \label{eq:lambert-phase}
\end{equation}
Equivalently,
\[
 \lambda(x)=-\frac1L W_{-1}(-Lx),
\]
but \eqref{eq:lambert-phase} is the identity used in the proofs.

The sharp expansion established for $F$ has the form
\begin{equation}
 \log F(x)
 =-\frac L2\lambda^2+
   \left(1+\frac L2\right)\lambda
   -\frac12\log\lambda
   +C_F+\Psi(\lambda)
   +\bigO(\lambda^{-1}), \label{eq:sharp-input}
\end{equation}
where $\Psi$ is smooth, bounded, and one-periodic.  The exact value of
$C_F$ and the Fourier description of $\Psi$ are important for $F$ itself,
but surprisingly they do not affect the leading equivalent of the
inverse.

\subsection{Asymptotic solution for the phase}

Let $y=F(x)$ and
\[
 T=\log\frac1y.
\]
Negating \eqref{eq:sharp-input} gives
\begin{equation}
 T=\frac L2\lambda^2-B\lambda+\frac12\log\lambda
     -C_F-\Psi(\lambda)+\bigO(\lambda^{-1}),
 \qquad B=1+\frac L2. \label{eq:T-lambda}
\end{equation}

\begin{lemma}[Asymptotic inversion of the quadratic phase; \Derived]
\label{lem:lambda-from-T}
As $T\to+\infty$,
\[
 \lambda
 =\sqrt{\frac{2T}{L}}+\frac BL
   +\bigO\!\left(\frac{\log T}{\sqrt T}\right).
\]
\end{lemma}

\begin{proof}
Since $\Psi$ is bounded, \eqref{eq:T-lambda} first gives
\[
 T=\frac L2\lambda^2+\bigO(\lambda+\log\lambda).
\]
Hence $\lambda\to\infty$ and, with
\[
 q=\sqrt{\frac{2T}{L}},
\]
we have $\lambda/q\to1$.  Dividing
\[
 T-\frac L2\lambda^2
 =-\frac L2(\lambda-q)(\lambda+q)
\]
by $\lambda+q\asymp q$ and using the preceding error estimate improves
this to $\lambda=q+\bigO(1)$.

Write
\[
 \lambda=q+\frac BL+\rho.
\]
Insert this into \eqref{eq:T-lambda}.  The terms proportional to $q$
cancel because $L(B/L)q=Bq$.  What remains is
\[
 Lq\rho
 =-\frac12\log\!\left(q+\frac BL+\rho\right)+\bigO(1)
   +\bigO(\rho^2).
\]
The already known bound $\rho=\bigO(1)$ makes the final term bounded.
Since $q\asymp\sqrt T$ and $\log q\asymp\log T$, division by $Lq$ gives
\[
 \rho=\bigO\!\left(\frac{\log T}{\sqrt T}\right).
\]
\end{proof}

\subsection{The endpoint equivalent for the inverse}

\begin{theorem}[Sharp inverse asymptotic; \Derived]
\label{thm:sharp-inverse-asymptotic}
As $y\downarrow0$,
\begin{equation}
 \boxed{\;
 \InvF(y)\sim
 \frac{\sqrt{\log(1/y)}}{\e\sqrt{\log2}}\,
 \exp\!\left[-\sqrt{2\log2\,\log(1/y)}\right].
 \;} \label{eq:sharp-inverse}
\end{equation}
Equivalently,
\begin{equation}
 \log\InvF(y)
 =-\sqrt{2L\log(1/y)}
  +\frac12\log\log(1/y)
  -1-\frac12\log L+o(1). \label{eq:log-sharp-inverse}
\end{equation}
At the right endpoint,
\begin{equation}
 1-\InvF(y)\sim
 \frac{\sqrt{\log(1/(1-y))}}{\e\sqrt{\log2}}\,
 \exp\!\left[-\sqrt{2\log2\,\log(1/(1-y))}\right]
 \quad(y\uparrow1). \label{eq:right-inverse-asymptotic}
\end{equation}
\end{theorem}

\begin{proof}
Let $x=\InvF(y)$.  Then $y=F(x)$ and $x$ is related to its exact phase by
$x=\lambda\e^{-L\lambda}$.  Therefore
\[
 \log x=\log\lambda-L\lambda.
\]
With $T=\log(1/y)$, \cref{lem:lambda-from-T} gives
\[
 \lambda=q+\frac BL+o(1),\qquad
 q=\sqrt{\frac{2T}{L}},
\]
because $(\log T)/\sqrt T=o(1)$.  Also $\log\lambda=\log q+o(1)$.
Consequently
\begin{align*}
 \log x
 &=\log q-Lq-B+o(1)\\
 &=-\sqrt{2LT}
   +\frac12\log T+\frac12\log\frac2L
   -1-\frac L2+o(1)\\
 &=-\sqrt{2LT}
   +\frac12\log T-1-\frac12\log L+o(1),
\end{align*}
because $L=\log2$.  This is \eqref{eq:log-sharp-inverse}.
Exponentiating proves \eqref{eq:sharp-inverse}.  Finally,
\eqref{eq:inv-reflection} gives
$1-\InvF(y)=\InvF(1-y)$ and hence
\eqref{eq:right-inverse-asymptotic}.
\end{proof}

\begin{remark}[Why the periodic fluctuation disappears at leading order]
The bounded term $\Psi(\lambda)$ changes the solution $\lambda$ of
\eqref{eq:T-lambda} only by $\bigO(1/\lambda)$.  Since
$\log x=\log\lambda-L\lambda$, this changes $\log x$ by $o(1)$ and
therefore changes $x$ by a factor tending to $1$.  Periodicity returns in
the next relative correction, but not in the leading equivalent
\eqref{eq:sharp-inverse}.
\end{remark}

\subsection{The complete hierarchy of elementary scales}

\begin{corollary}[Between powers and logarithms; \Derived]
\label{cor:inverse-scale-hierarchy}
For every $\alpha>0$ and every $m>0$,
\begin{align}
 y^\alpha&=o(\InvF(y)), \label{eq:power-below-G}\\
 \InvF(y)&=o\!\left((\log(1/y))^{-m}\right)
 \label{eq:G-below-log}
\end{align}
as $y\downarrow0$.
\end{corollary}

\begin{proof}
Set $T=\log(1/y)$.  Formula \eqref{eq:log-sharp-inverse} gives
\[
 \log\frac{\InvF(y)}{y^\alpha}
 =\alpha T-\sqrt{2LT}+\bigO(\log T)\longrightarrow+\infty,
\]
which proves \eqref{eq:power-below-G}.  Likewise,
\[
 \log\!\left(\InvF(y)T^m\right)
 =-\sqrt{2LT}+\bigO(\log T)\longrightarrow-\infty,
\]
which proves \eqref{eq:G-below-log}.
\end{proof}

\begin{boxedremark}
The inverse endpoint is therefore neither power-like nor logarithmic.
It lies in the stretched-logarithmic regime
$\exp(-c\sqrt{\log(1/y)})$ with an essential
$\sqrt{\log(1/y)}$ prefactor.
\end{boxedremark}

\section{Closed periodic saddle jets}
\label{sec:closed-jets}

We now turn from the inverse to the coefficient machinery hidden inside
the all-orders endpoint expansion.  Throughout this part,
\[
 L=\log2,\qquad D=\frac{\dd}{\dd u},
\]
and $\Psi$ is the smooth one-periodic function occurring in the sharp
Lambert main term.

\subsection{The recursive jets}

Define the periodic jets $p_n$ by
\begin{align}
 p_0(u)&=\frac12+\frac{\Psi'(u)}L, \label{eq:p0}\\
 p_{n+1}(u)
 &=\left(\frac DL-(n+1)\right)p_n(u)
   +\frac{(-1)^{n+1}n!}{L}. \label{eq:p-recurrence}
\end{align}
The bounded exponent jets are
\begin{equation}
 d_n(u)=p_n(u)+(-1)^{n+1}n!. \label{eq:d-from-p}
\end{equation}
The subtraction in \eqref{eq:d-from-p} removes the factorial part that
would otherwise obscure the centered saddle expansion.

For $n\geq0$, write
\[
 H_n=\sum_{k=1}^n\frac1k,\qquad H_0=0,
\]
and define integers $c(n,m)$ by
\begin{equation}
 \prod_{k=1}^n(z-k)=\sum_{m=0}^n c(n,m)z^m.
 \label{eq:c-definition}
\end{equation}

\begin{theorem}[Closed periodic jets; \Formalized]
\label{thm:closed-jets}
For every $n\geq0$,
\begin{align}
 p_n(u)
 &=
 (-1)^n n!\left(\frac12+\frac{H_n}{L}\right)
 +\sum_{m=0}^n
   c(n,m)\frac{\Psi^{(m+1)}(u)}{L^{m+1}},
 \label{eq:p-closed}\\
 d_n(u)
 &=
 (-1)^n n!\left(\frac{H_n}{L}-\frac12\right)
 +\sum_{m=0}^n
   c(n,m)\frac{\Psi^{(m+1)}(u)}{L^{m+1}}.
 \label{eq:d-closed}
\end{align}
\end{theorem}

\begin{proof}
Let
\[
 T_k=\frac DL-k.
\]
The operators $T_k$ commute because each is a polynomial in $D$.  Iterating
\eqref{eq:p-recurrence} gives
\begin{equation}
 p_n=T_nT_{n-1}\cdots T_1p_0
 +\sum_{j=0}^{n-1}
   T_nT_{n-1}\cdots T_{j+2}
   \left(\frac{(-1)^{j+1}j!}{L}\right),
 \label{eq:p-iteration}
\end{equation}
where the product is empty when $j=n-1$.

First consider the homogeneous term.  Since
\[
 T_n\cdots T_1
 =\prod_{k=1}^n\left(\frac DL-k\right),
\]
its action on the constant $1/2$ is
\[
 \frac12\prod_{k=1}^n(-k)=\frac{(-1)^nn!}{2}.
\]
Its action on $\Psi'/L$ is found by expanding the operator polynomial with
\eqref{eq:c-definition}:
\[
 \prod_{k=1}^n\left(\frac DL-k\right)\frac{\Psi'}L
 =\sum_{m=0}^n
   c(n,m)\frac{\Psi^{(m+1)}}{L^{m+1}}.
\]

Now propagate the forcing term inserted at stage $j$.  It is constant, so
each subsequent $T_k$ acts on it as multiplication by $-k$.  Hence its
contribution at stage $n$ is
\begin{align*}
 \frac{(-1)^{j+1}j!}{L}
 \prod_{k=j+2}^n(-k)
 &=
 \frac{(-1)^{j+1}j!}{L}
 \frac{(-1)^{n-j-1}n!}{(j+1)!}\\
 &=\frac{(-1)^nn!}{(j+1)L}.
\end{align*}
Summing over $j=0,\dots,n-1$ gives
\[
 \frac{(-1)^nn!}{L}\sum_{j=0}^{n-1}\frac1{j+1}
 =\frac{(-1)^nn!H_n}{L}.
\]
Combining the constant, derivative, and forcing contributions proves
\eqref{eq:p-closed}.  Finally,
\[
 p_n+(-1)^{n+1}n!
 =(-1)^nn!\left(\frac12+\frac{H_n}{L}-1\right)
   +\text{derivative part},
\]
which is exactly \eqref{eq:d-closed}.
\end{proof}

\subsection{Why the Stirling indices are shifted}

Let $s(n,k)$ be the signed Stirling number of the first kind, defined by
\[
 z(z-1)\cdots(z-n+1)=\sum_{k=0}^n s(n,k)z^k.
\]

\begin{proposition}[Shifted Stirling conversion; \Formalized]
\label{prop:stirling-shift}
For $0\leq m\leq n$,
\[
 c(n,m)=s(n+1,m+1).
\]
\end{proposition}

\begin{proof}
Apply the defining identity with $n+1$:
\[
 z(z-1)\cdots(z-n)
 =\sum_{k=0}^{n+1}s(n+1,k)z^k.
\]
Both sides have zero constant term, so divide by $z$:
\[
 (z-1)\cdots(z-n)
 =\sum_{m=0}^n s(n+1,m+1)z^m.
\]
Comparison with \eqref{eq:c-definition} proves the formula.
\end{proof}

\begin{warning}
The coefficients are not $s(n,m)$.  The product starts with $(z-1)$
rather than $z$, and this forces the simultaneous shift
$(n,m)\mapsto(n+1,m+1)$.  In particular, the nonzero constant column
$c(n,0)=(-1)^nn!$ is essential.
\end{warning}

\subsection{Coefficient rows and the first jets}

The first rows of $c(n,m)$, listed from $m=0$ to $m=n$, are
\[
\begin{array}{c|rrrrrr}
n&c(n,0)&c(n,1)&c(n,2)&c(n,3)&c(n,4)&c(n,5)\\
\hline
0&1\\
1&-1&1\\
2&2&-3&1\\
3&-6&11&-6&1\\
4&24&-50&35&-10&1\\
5&-120&274&-225&85&-15&1
\end{array}
\]
and \eqref{eq:d-closed} gives
\begin{align}
 d_0={}&-\frac12+\frac{\Psi'}L, \label{eq:d0-explicit}\\
 d_1={}&\frac12-\frac1L-\frac{\Psi'}L+\frac{\Psi''}{L^2},
 \label{eq:d1-explicit}\\
 d_2={}&-1+\frac3L+\frac{2\Psi'}L
       -\frac{3\Psi''}{L^2}+\frac{\Psi'''}{L^3},
 \label{eq:d2-explicit}\\
 d_3={}&3-\frac{11}{L}-\frac{6\Psi'}L
       +\frac{11\Psi''}{L^2}-\frac{6\Psi'''}{L^3}
       +\frac{\Psi^{(4)}}{L^4}. \label{eq:d3-explicit}
\end{align}

\begin{corollary}[Structural properties; \Formalized]
\label{cor:jet-structure}
Every $p_n$ and $d_n$ is smooth, one-periodic, and bounded on $\R$.
Moreover, $d_n$ involves derivatives of $\Psi$ only through order $n+1$,
and the coefficient of the highest derivative is $L^{-(n+1)}$.
\end{corollary}

\begin{proof}
The closed formulas are finite linear combinations of derivatives of the
smooth periodic function $\Psi$ and constants.  Smoothness, periodicity,
and boundedness follow.  Since $c(n,n)=1$, the coefficient of
$\Psi^{(n+1)}$ is $L^{-(n+1)}$.
\end{proof}


\section{The centered exponent and its algebraic collapse}
\label{sec:centered-exponent}

\subsection{From a long Taylor expansion to one short formula}

At the saddle, introduce the half-power parameter
\[
 \varepsilon=\lambda^{-1/2}
\]
and the Gaussian variable $v$.  After removing the exact quadratic term,
the centered exponent has a formal expansion
\[
 \mathcal E(u,v;\varepsilon)
 =\sum_{m\geq1}E_m(u,v)\varepsilon^m.
\]
The original construction of $E_m$ contains separate contributions from
the vertical negative-Laplace jet and from the Taylor expansion of
$\log(1+\ii\varepsilon v)$.  The bounded exponent jets make these two
pieces collapse.

\begin{theorem}[Closed exponent coefficient; \Formalized]
\label{thm:closed-exponent}
For every $m\geq1$,
\begin{equation}
 E_m(u,v)=\ii^m\left(
   \frac{d_{m-1}(u)}{m!}v^m
   +\frac{(-1)^{m+1}}{m+2}v^{m+2}
 \right). \label{eq:E-closed}
\end{equation}
In particular,
\begin{align}
 E_1&=\ii v\,\frac{3d_0+v^2}{3}, \label{eq:E1}\\
 E_2&=v^2\,\frac{-2d_1+v^2}{4}, \label{eq:E2}\\
 E_3&=\ii v^3\left(-\frac{d_2}{6}-\frac{v^2}{5}\right), \label{eq:E3}\\
 E_4&=\frac{v^4}{24}(d_3-4v^2). \label{eq:E4}
\end{align}
\end{theorem}

\begin{proof}
Before centering, the coefficient at order $m$ has two contributions.
The vertical jet contributes $\ii^m p_{m-1}(u)v^m/m!$, while the
logarithmic saddle coordinate contributes a factorial term to the same
monomial and a universal $v^{m+2}$ term.  Using
\[
 d_{m-1}=p_{m-1}+(-1)^m(m-1)!
\]
combines the two $v^m$ coefficients into
$\ii^m d_{m-1}/m!$.  The remaining logarithmic term is
\[
 \frac{\ii^{m+2}(-1)^m}{m+2}v^{m+2}
 =\ii^m\frac{(-1)^{m+1}}{m+2}v^{m+2},
\]
because $\ii^2=-1$.  This is \eqref{eq:E-closed}.  Substitution of
$m=1,2,3,4$ gives \eqref{eq:E1}--\eqref{eq:E4}.
\end{proof}

\begin{corollary}[Parity; \Formalized]
\label{cor:E-parity}
For every $m\geq1$,
\[
 E_m(u,-v)=(-1)^mE_m(u,v).
\]
Thus the odd-indexed $E_m$ are odd polynomials in $v$ and the
even-indexed $E_m$ are even.
\end{corollary}

\begin{proof}
Both monomials in \eqref{eq:E-closed} have parity $m$, because $m+2$ has
the same parity as $m$.
\end{proof}

\section{Formal exponential, Gaussian contraction, and logarithm}
\label{sec:coefficient-formulas}

\subsection{The exponential recursion}

Define polynomials $B_n(u,v)$ by the formal identity
\begin{equation}
 \exp\!\left(\sum_{m\geq1}E_m(u,v)\varepsilon^m\right)
 =\sum_{n\geq0}B_n(u,v)\varepsilon^n.
 \label{eq:B-definition}
\end{equation}
Thus $B_0=1$.

\begin{proposition}[Finite exponential recursion; \Formalized]
\label{prop:exp-recursion}
For $n\geq1$,
\begin{equation}
 nB_n=\sum_{m=1}^n mE_mB_{n-m}. \label{eq:B-recursion}
\end{equation}
Equivalently,
\[
 B_n=\sum_{\substack{r\geq1\\m_1+\cdots+m_r=n}}
 \frac1{r!}E_{m_1}\cdots E_{m_r},
\]
where the inner sum is over ordered compositions of $n$ into $r$ positive
parts.
\end{proposition}

\begin{proof}
Differentiate \eqref{eq:B-definition} with respect to $\varepsilon$.
The derivative of the left side is
\[
 \left(\sum_{m\geq1}mE_m\varepsilon^{m-1}\right)
 \left(\sum_{k\geq0}B_k\varepsilon^k\right),
\]
whereas the derivative of the right side is
$\sum_{n\geq1}nB_n\varepsilon^{n-1}$.  Comparing the coefficient of
$\varepsilon^{n-1}$ gives \eqref{eq:B-recursion}.  Expanding the
exponential as
\[
 \sum_{r\geq0}\frac1{r!}
 \left(\sum_{m\geq1}E_m\varepsilon^m\right)^r
\]
and collecting total degree $n$ gives the composition formula.
\end{proof}

\begin{corollary}[Parity of the exponential coefficients; \Formalized]
\label{cor:B-parity}
For all $n\geq0$,
\[
 B_n(u,-v)=(-1)^nB_n(u,v).
\]
\end{corollary}

\begin{proof}
Every product contributing to $B_n$ has indices
$m_1+\cdots+m_r=n$.  By \cref{cor:E-parity}, its parity is
$(-1)^{m_1+\cdots+m_r}=(-1)^n$.
\end{proof}

\subsection{Gaussian contraction and mass coefficients}

For a polynomial $P$ define normalized Gaussian contraction by
\[
 \Gauss[P]
 =\frac1{\sqrt{2\pi}}\int_{\R}\e^{-v^2/2}P(v)\,\dd v.
\]
It is determined by
\begin{equation}
 \Gauss[v^{2k}]=(2k-1)!!,\qquad
 \Gauss[v^{2k+1}]=0. \label{eq:gaussian-moments}
\end{equation}
The mass coefficients are
\begin{equation}
 a_j(u)=\Gauss[B_{2j}(u,\cdot)],\qquad j\geq0.
 \label{eq:mass-coeff-definition}
\end{equation}
In particular $a_0=1$.  The odd half-powers vanish automatically:
\[
 \Gauss[B_{2j+1}]=0
\]
by \cref{cor:B-parity}.

\begin{theorem}[Closed composition formula for \(a_j\); \Derived]
\label{thm:mass-composition}
For $j\geq1$,
\begin{align}
 a_j(u)
 ={}&(-1)^j
 \sum_{r=1}^{2j}\frac1{r!}
 \sum_{\substack{m_1+\cdots+m_r=2j\\m_i\geq1}}
 \sum_{S\subseteq\{1,\dots,r\}}
 (2j+2|S|-1)!! \notag\\
 &\qquad{}\times
 \prod_{i\in S}\frac{(-1)^{m_i+1}}{m_i+2}
 \prod_{i\notin S}\frac{d_{m_i-1}(u)}{m_i!}.
 \label{eq:mass-composition}
\end{align}
Consequently $a_j$ is a polynomial with rational coefficients in
$d_0,\dots,d_{2j-1}$.
\end{theorem}

\begin{proof}
Use the ordered-composition formula in
\cref{prop:exp-recursion} for $B_{2j}$.  Each factor
\[
 E_{m_i}
 =\ii^{m_i}\left(
   \frac{d_{m_i-1}}{m_i!}v^{m_i}
   +\frac{(-1)^{m_i+1}}{m_i+2}v^{m_i+2}
 \right)
\]
offers two choices.  Let $S$ be the set of factors from which the second,
universal monomial is chosen.  Since $\sum_i m_i=2j$, the product of the
powers of $\ii$ is
\[
 \ii^{2j}=(-1)^j,
\]
and the total power of $v$ is $2j+2|S|$.  Its Gaussian contraction is
$(2j+2|S|-1)!!$.  Multiplying the remaining scalar factors gives
\eqref{eq:mass-composition}.  The largest possible jet index is
$m_i-1\leq2j-1$.
\end{proof}

\begin{boxedremark}
Formula \eqref{eq:mass-composition} is a human-readable algebraic
consequence of the formal exponent and exponential-recursion theorems.
The repository machine-checks the underlying recursion, the structural
top-jet theorem, and explicit low orders.  It does not package this exact
threefold summation as a single general Lean declaration.
\end{boxedremark}

\subsection{Logarithmic coefficients}

Define $A_j(u)$ by
\begin{equation}
 1+\sum_{j\geq1}a_j(u)z^j
 =\exp\!\left(\sum_{j\geq1}A_j(u)z^j\right).
 \label{eq:log-coeff-definition}
\end{equation}
These are the coefficients that occur in $\log F$.

\begin{proposition}[Mass-to-logarithm conversion; \Formalized]
\label{prop:log-recursion}
For $n\geq1$,
\begin{equation}
 A_n
 =a_n-\frac1n\sum_{k=1}^{n-1}kA_k\,a_{n-k}.
 \label{eq:A-recursion}
\end{equation}
Equivalently,
\begin{equation}
 A_n
 =\sum_{r=1}^n\frac{(-1)^{r-1}}r
   \sum_{\substack{q_1+\cdots+q_r=n\\q_i\geq1}}
   a_{q_1}\cdots a_{q_r}. \label{eq:A-compositions}
\end{equation}
\end{proposition}

\begin{proof}
Differentiate \eqref{eq:log-coeff-definition} logarithmically:
\[
 \sum_{n\geq1}na_nz^{n-1}
 =
 \left(1+\sum_{j\geq1}a_jz^j\right)
 \left(\sum_{k\geq1}kA_kz^{k-1}\right).
\]
The coefficient of $z^{n-1}$ gives
\[
 na_n=\sum_{k=1}^n kA_ka_{n-k},
 \qquad a_0=1,
\]
and solving for $A_n$ proves \eqref{eq:A-recursion}.  Formula
\eqref{eq:A-compositions} is the coefficient of $z^n$ in the ordinary
series
\[
 \log(1+w)=\sum_{r\geq1}\frac{(-1)^{r-1}}r w^r,
 \qquad w=\sum_{j\geq1}a_jz^j.
\]
\end{proof}

\begin{corollary}[Regularity and derivative order; \Formalized\ input,
\Derived\ conclusion]
\label{cor:A-structure}
For every $j$, the functions $a_j$ and $A_j$ are smooth, one-periodic,
and bounded.  Neither involves a derivative of $\Psi$ of order greater
than $2j$.
\end{corollary}

\begin{proof}
The function $a_j$ is a polynomial in
$d_0,\dots,d_{2j-1}$, and $d_n$ is smooth, periodic, and depends on
$\Psi$ only through order $n+1$.  Thus $a_j$ has the stated properties
and uses derivatives only through order $2j$.  The recursion
\eqref{eq:A-recursion} transfers the same assertions to $A_j$.
\end{proof}

\subsection{The top jet and the highest derivative}

The most complicated-looking coefficient has a very simple dependence on
its newest jet.

\begin{theorem}[Leading jet law; \Formalized]
\label{thm:top-jet}
For $j\geq1$, the mass coefficient $a_j$ is affine in $d_{2j-1}$, and
the coefficient of that jet is
\begin{equation}
 \frac{(-1)^j}{2^j j!}. \label{eq:top-jet-coefficient}
\end{equation}
The logarithmic coefficient $A_j$ has the same top-jet coefficient.
Consequently the coefficient of $\Psi^{(2j)}(u)$ in $A_j(u)$ is
\begin{equation}
 \boxed{\frac{(-1)^j}{2^j j!L^{2j}}.}
 \label{eq:highest-derivative}
\end{equation}
\end{theorem}

\begin{proof}
The jet $d_{2j-1}$ occurs only in $E_{2j}$.  Therefore the only term of
$B_{2j}$ that contains it is the linear term $E_{2j}$ itself.  From
\eqref{eq:E-closed}, its jet part is
\[
 \ii^{2j}\frac{d_{2j-1}}{(2j)!}v^{2j}
 =(-1)^j\frac{d_{2j-1}}{(2j)!}v^{2j}.
\]
Gaussian contraction multiplies this by $(2j-1)!!$.  Since
\[
 (2j)!=2^jj!(2j-1)!!,
\]
the coefficient is \eqref{eq:top-jet-coefficient}.

In \eqref{eq:A-recursion}, every correction term
$A_ka_{j-k}$ has $k<j$ and $j-k<j$; it involves jets only through
$d_{2j-3}$ and therefore cannot alter the coefficient of $d_{2j-1}$.
Thus $A_j$ has the same top-jet coefficient.  Finally,
\cref{cor:jet-structure} says that the coefficient of
$\Psi^{(2j)}$ in $d_{2j-1}$ is $L^{-2j}$, giving
\eqref{eq:highest-derivative}.
\end{proof}

\begin{remark}
The theorem provides a quick consistency check for every explicit
coefficient.  For example, $A_1$ must contain
$-\Psi''/(2L^2)$, while $A_2$ must contain
$+\Psi^{(4)}/(8L^4)$.
\end{remark}

\section{The first and second logarithmic corrections}
\label{sec:second-correction}

\subsection{Order two: the first correction}

From \cref{eq:E1,eq:E2},
\[
 B_2=E_2+\frac12E_1^2.
\]
A direct expansion gives
\[
 B_2
 =-\frac{d_0^2+d_1}{2}v^2
   +\left(\frac14-\frac{d_0}{3}\right)v^4
   -\frac1{18}v^6.
\]
Using $\Gauss[v^2]=1$, $\Gauss[v^4]=3$, and
$\Gauss[v^6]=15$ yields
\begin{equation}
 a_1=A_1
 =-\frac{d_0^2}{2}-d_0-\frac{d_1}{2}-\frac1{12}.
 \label{eq:A1-jets}
\end{equation}
Equivalently,
\[
 A_1=-\frac{6d_0^2+12d_0+6d_1+1}{12}.
\]
Substituting \cref{eq:d0-explicit,eq:d1-explicit} recovers the first
periodic correction in terms of $\Psi'$ and $\Psi''$.

\subsection{Order four of the formal exponential}

The fourth exponential coefficient is
\begin{equation}
 B_4
 =E_4+E_3E_1+\frac12E_2^2
   +\frac12E_2E_1^2+\frac1{24}E_1^4.
 \label{eq:B4-identity}
\end{equation}
This follows either from the composition formula or from the recurrence
\eqref{eq:B-recursion}.  Substitution of
\cref{eq:E1,eq:E2,eq:E3,eq:E4} and collection of even powers of $v$
gives the following polynomial.

\begin{proposition}[The order-four reference polynomial; \Formalized]
\label{prop:B4-polynomial}
Writing $d_k=d_k(u)$,
\begin{align}
 B_4(u,v)={}&
 \left(\frac{d_0^4}{24}+\frac{d_0^2d_1}{4}
       +\frac{d_0d_2}{6}+\frac{d_1^2}{8}
       +\frac{d_3}{24}\right)v^4 \notag\\
 &+\left(\frac{d_0^3}{18}-\frac{d_0^2}{8}
       +\frac{d_0d_1}{6}+\frac{d_0}{5}
       -\frac{d_1}{8}+\frac{d_2}{18}
       -\frac16\right)v^6 \notag\\
 &+\left(\frac{d_0^2}{36}-\frac{d_0}{12}
       +\frac{d_1}{36}+\frac{47}{480}\right)v^8 \notag\\
 &+\left(\frac{d_0}{162}-\frac1{72}\right)v^{10}
   +\frac1{1944}v^{12}. \label{eq:B4-polynomial}
\end{align}
\end{proposition}

\begin{proof}
Let
\[
 P_1=d_0v+\frac{v^3}{3},\qquad
 P_3=-\frac{d_2v^3}{6}-\frac{v^5}{5},
\]
so that $E_1=\ii P_1$ and $E_3=\ii P_3$.  Then
$E_3E_1=-P_3P_1$ and $E_1^4=P_1^4$.  The even terms $E_2,E_4$ are already
real.  Insert these expressions in \eqref{eq:B4-identity}.  Each summand
has degree at most $12$ and is even.  Multiplication gives
\[
\begin{array}{c|ccccc}
 &v^4&v^6&v^8&v^{10}&v^{12}\\
\hline
E_4&
 d_3/24&-1/6&0&0&0\\
E_3E_1&
 d_0d_2/6&d_2/18+d_0/5&1/15&0&0\\
E_2^2/2&
 d_1^2/8&-d_1/8&1/32&0&0\\
E_2E_1^2/2&
 d_0^2d_1/4&
 d_0d_1/6-d_0^2/8&
 d_1/36-d_0/12&
 -1/72&0\\
E_1^4/24&
 d_0^4/24&d_0^3/18&d_0^2/36&d_0/162&1/1944
\end{array}
\]
and the two universal $v^8$ contributions satisfy
$1/15+1/32=47/480$.  Summing the columns proves
\eqref{eq:B4-polynomial}.
\end{proof}

\subsection{Gaussian contraction and the second mass coefficient}

The relevant Gaussian moments are
\[
 \Gauss[v^4]=3,\quad
 \Gauss[v^6]=15,\quad
 \Gauss[v^8]=105,\quad
 \Gauss[v^{10}]=945,\quad
 \Gauss[v^{12}]=10395.
\]

\begin{theorem}[Second mass coefficient; \Formalized]
\label{thm:a2}
Gaussian contraction of \eqref{eq:B4-polynomial} gives
\begin{align}
 a_2={}&
 \frac{d_0^4}{8}+\frac{5d_0^3}{6}
 +\frac{3d_0^2d_1}{4}+\frac{25d_0^2}{24}
 +\frac{5d_0d_1}{2}+\frac{d_0d_2}{2}
 +\frac{d_0}{12} \notag\\
 &+\frac{3d_1^2}{8}+\frac{25d_1}{24}
 +\frac{5d_2}{6}+\frac{d_3}{8}
 +\frac1{288}. \label{eq:a2}
\end{align}
\end{theorem}

\begin{proof}
Multiply the coefficients of $v^4,v^6,v^8,v^{10},v^{12}$ in
\eqref{eq:B4-polynomial} by $3,15,105,945,10395$, respectively, and
collect like monomials.  For illustration, the coefficient of $d_0$ is
\[
 15\cdot\frac15
 +105\left(-\frac1{12}\right)
 +945\cdot\frac1{162}
 =3-\frac{35}{4}+\frac{35}{6}
 =\frac1{12},
\]
and the universal constant is
\[
 15\left(-\frac16\right)
 +105\frac{47}{480}
 +945\left(-\frac1{72}\right)
 +10395\frac1{1944}
 =\frac1{288}.
\]
The remaining coefficients follow by the same arithmetic.
\end{proof}

\subsection{The second logarithmic coefficient}

Since
\[
 \log(1+a_1z+a_2z^2+\cdots)
 =a_1z+\left(a_2-\frac12a_1^2\right)z^2+\cdots,
\]
we have $A_2=a_2-a_1^2/2$.

\begin{theorem}[Second periodic saddle correction; \Formalized]
\label{thm:A2}
The second logarithmic coefficient is
\begin{equation}
 \boxed{
 A_2=
 \frac{
  8d_0^3+12d_0^2d_1+12d_0^2+48d_0d_1+12d_0d_2
  +6d_1^2+24d_1+20d_2+3d_3
 }{24}.}
 \label{eq:A2}
\end{equation}
It is smooth, one-periodic, and bounded.
\end{theorem}

\begin{proof}
Insert \eqref{eq:A1-jets} and \eqref{eq:a2} into
$A_2=a_2-a_1^2/2$.  The quartic term $d_0^4/8$ cancels, as do the
terms $d_0/12$ and $1/288$ after the square is expanded.  The remaining
terms reduce to \eqref{eq:A2}.  Smoothness, periodicity, and boundedness
follow because $A_2$ is a polynomial in $d_0,d_1,d_2,d_3$.
\end{proof}

\begin{corollary}[Two explicit correction orders; \Formalized]
\label{cor:two-corrections}
Let $\lambda=\lambda(x)$ be the exact lower Lambert phase.  Then, as
$x\downarrow0$,
\begin{equation}
 \log F(x)
 =\Main(x)+\frac{A_1(\lambda)}{\lambda}
          +\frac{A_2(\lambda)}{\lambda^2}
          +\bigO(\lambda^{-3}), \label{eq:two-correction-expansion}
\end{equation}
where $\Main(x)$ is the exact sharp Lambert main term.  Since
$\lambda\asymp-\log x$, the remainder is also
$\bigO((-\log x)^{-3})$.
\end{corollary}

\begin{warning}
The coefficients in \eqref{eq:two-correction-expansion} are evaluated at
the exact phase $\lambda(x)$.  Substituting a truncated asymptotic formula
for $\lambda$ into $A_1$ or $A_2$ is not licensed merely by
$\lambda-\lambda_{\rm approx}=\bigO(1)$: periodic functions are sensitive
to bounded phase shifts.
\end{warning}


\section{Why the all-orders remainder is actually rigorous}
\label{sec:all-orders-proof}

The formal coefficient algebra by itself does not prove an asymptotic
expansion.  One must control the exact contour integral uniformly on a
window that grows with the requested order, prove that exponentiation does
not magnify the Taylor error, replace the finite Gaussian window by the
whole line, and show that the discarded contour is smaller than every
retained inverse power.  This section spells out that analytic chain.

\subsection{The exact normalized saddle mass}

For $x>0$ small, let $\lambda=\lambda(x)$ be the exact phase
\eqref{eq:lambert-phase}.  The exact saddle reduction can be organized as
\begin{equation}
 \log F(x)
 =\Main(x)+\log\mathcal G(x)+\mathcal T(x),
 \label{eq:exact-saddle-reduction}
\end{equation}
where:
\begin{itemize}
\item $\Main(x)$ is the exact sharp Lambert main term, containing the
quadratic phase, the Gaussian normalization, the periodic finite part, and
the fixed constant;
\item $\mathcal G(x)>0$ is the normalized saddle-kernel mass and
$\mathcal G(x)\to1$;
\item $\mathcal T(x)$ is the forward negative-Laplace tail, which is flat
against every inverse power of $\lambda$.
\end{itemize}
Thus all algebraic corrections beyond $\Main$ come from $\log\mathcal G$.

It is technically convenient to begin on the real dyadic logarithmic
scale
\[
 x=2^{-t},\qquad t\to+\infty,
\]
and to write
\[
 b=\lambda(2^{-t}),\qquad \varepsilon=b^{-1/2}.
\]
This is not a restriction to integer $t$.  Every sufficiently small
positive $x$ has the unique representation $x=2^{-t}$ with
$t=-\log_2x$.

After the saddle contour is centered and rescaled by
$v=\sqrt b\,\theta$, its central kernel has the normalized form
\begin{equation}
 \frac1{\sqrt{2\pi}}\e^{-v^2/2}
 \exp R_b(v), \label{eq:normalized-central-kernel}
\end{equation}
where $R_b$ is the exact centered exponent.  The formal expansion of
$R_b$ is
\begin{equation}
 R_b(v)\sim\sum_{m\geq1}E_m(b,v)b^{-m/2}.
 \label{eq:R-formal}
\end{equation}
The first variable of $E_m$ is evaluated at $b$ because all coefficient
functions are periodic in the exact phase.

\subsection{The order-dependent central window}

For a requested order $N\geq0$, define
\begin{equation}
 A_N(b)=\sqrt{32(N+1)\log b}. \label{eq:central-radius}
\end{equation}

\begin{lemma}[Scale separation; \Formalized]
\label{lem:central-scale}
For every fixed $N$,
\[
 A_N(b)\longrightarrow+\infty,\qquad
 \varepsilon A_N(b)\longrightarrow0
 \qquad(b\to+\infty).
\]
Moreover, for every fixed $r>0$ and $\delta>0$,
\[
 A_N(b)^r=o(b^\delta).
\]
\end{lemma}

\begin{proof}
The first limit is immediate from $\log b\to\infty$.  The second follows
from
\[
 \varepsilon A_N(b)
 =\sqrt{\frac{32(N+1)\log b}{b}}\longrightarrow0.
\]
Finally, every fixed power of $\log b$ grows more slowly than every
positive power of $b$.
\end{proof}

The window has two simultaneous purposes.  It grows, so replacing a
Gaussian integral over it by a whole-line Gaussian integral creates a
tiny tail.  Yet its angular width
$A_N(b)/\sqrt b$ shrinks to zero, so all local Taylor estimates remain
uniform.

\subsection{All-order control of the exact centered exponent}

The exact negative-Laplace factor contains a finite periodic part plus
endpoint and forward-product corrections.  On the vertical line through
the saddle, the proof decomposes it into:
\begin{enumerate}
\item a Taylor polynomial in $\theta=v/\sqrt b$ whose coefficients are
the periodic jets $p_n$;
\item the Taylor polynomial of $\log(1+\ii\theta)$;
\item two boundary-jet remainders from the finite part;
\item forward-product jets whose size is exponentially small in $b$;
\item explicit integral remainders for the vertical derivatives.
\end{enumerate}
The factorial pieces in the first two items combine into the bounded
jets $d_n$, producing \eqref{eq:E-closed}.

\begin{proposition}[Integrated exponent truncation; \Formalized]
\label{prop:exponent-truncation}
Fix $N$.  The Taylor depth can be chosen, depending only on $N$, so that
on $|v|\leq A_N(b)$ the exact exponent may be replaced by
\[
 R_{b,N}(v)
 =\sum_{m=1}^{2N+1}E_m(b,v)b^{-m/2}
\]
with an error which, after multiplication by the Gaussian weight and
integration over the central window, is
\[
 \bigO(b^{-N-1}).
\]
The implied constant may depend on $N$ but not on $b$.
\end{proposition}

\begin{proof}
All vertical derivatives of the periodic finite part are uniformly
bounded because the periodic jets are smooth and bounded.  Taylor's
theorem therefore bounds a vertical remainder of order $M$ by
\[
 C_M|\theta|^{M+1}
 =C_M b^{-(M+1)/2}|v|^{M+1}.
\]
The logarithmic coordinate has an analogous remainder on the shrinking
angular window, because $\varepsilon A_N(b)\to0$ keeps
$1+\ii\theta$ uniformly away from zero.  Boundary jets satisfy the same
kind of estimate, while the forward-product terms are exponentially
small and hence smaller than any prescribed inverse power.

Choose $M$ with enough surplus beyond $2N+1$.  On the central window,
powers of $|v|$ are bounded by powers of $A_N(b)$, hence by powers of
$\log b$.  By \cref{lem:central-scale}, the spare powers of $b^{-1/2}$
absorb every such logarithmic loss.  The pointwise exponent error is
therefore integrably $\bigO(b^{-N-1})$ against
$\e^{-v^2/2}\dd v$.  Removing terms above order $2N+1$ from the finite
Taylor polynomial costs the same order because each carries at least the
corresponding spare half-power.  This leaves $R_{b,N}$.
\end{proof}

\subsection{Finite exponentiation without an infinite-series argument}

A common informal step is to exponentiate \eqref{eq:R-formal} term by
term.  The formal development instead uses a finite Taylor polynomial for
the exponential and an explicit remainder.

Let
\[
 \mathcal P_N(b,v)
 =\sum_{\ell=0}^{2N+1}B_\ell(b,v)b^{-\ell/2}.
 \label{eq:reference-polynomial}
\]
It is the reference polynomial obtained from the formal recursion in
\cref{prop:exp-recursion}.

\begin{proposition}[Finite exponential substitution; \Formalized]
\label{prop:finite-exp-substitution}
For every fixed $N$,
\begin{equation}
 \frac1{\sqrt{2\pi}}
 \int_{|v|\leq A_N(b)}
 \e^{-v^2/2}
 \left|\exp R_b(v)-\mathcal P_N(b,v)\right|\dd v
 =\bigO(b^{-N-1}). \label{eq:central-reference-error}
\end{equation}
\end{proposition}

\begin{proof}
First replace $R_b$ by $R_{b,N}$ using
\cref{prop:exponent-truncation}.  On the central window the truncated
exponent tends uniformly to zero after the exact quadratic term has been
removed, and it is bounded by a fixed polynomial in $v$ times
$b^{-1/2}$.  Apply Taylor's formula
\[
 \e^z=\sum_{r=0}^{L-1}\frac{z^r}{r!}
      +\frac{z^L}{(L-1)!}\int_0^1(1-s)^{L-1}\e^{sz}\,\dd s
\]
with $L=2(N+1)$.  Uniform control of the real part of $R_{b,N}$ bounds the
integral remainder by a Gaussian-integrable polynomial times
$b^{-N-1}$.

Now expand the finite Taylor polynomial in the finitely many powers
$b^{-m/2}$.  The coefficient of $b^{-\ell/2}$ for
$\ell<L$ is exactly $B_\ell$, by the same recursion that defines the
formal exponential coefficients.  Algebraically, the difference between
the finite exponential Taylor polynomial and
$\mathcal P_N$ factors as
\[
 b^{-L/2}\times
 \{\text{a finite quotient polynomial in }b^{-1/2},v,
   E_1,\dots,E_{2N+1}\}.
\]
Because $L/2=N+1$ and the coefficient functions are bounded, Gaussian
integration gives \eqref{eq:central-reference-error}.
\end{proof}

\subsection{Gaussian contraction on the whole line}

By \cref{cor:B-parity}, odd $\ell$ contribute odd polynomials.  Therefore
\begin{equation}
 \Gauss[\mathcal P_N(b,\cdot)]
 =\sum_{j=0}^{N}a_j(b)b^{-j}. \label{eq:reference-contraction}
\end{equation}
The exact central integral, however, only runs over
$|v|\leq A_N(b)$.  We must show that the polynomial reference can be
extended to the whole real line at the required order.

\begin{lemma}[Gaussian polynomial tail; \Formalized]
\label{lem:gaussian-polynomial-tail}
For every fixed polynomial $P$ and every fixed $N$,
\[
 \int_{|v|>A_N(b)}\e^{-v^2/2}|P(v)|\,\dd v
 =\bigO(b^{-N-1}).
\]
The same assertion holds for a finite family of polynomials whose
coefficients are uniformly bounded in $b$.
\end{lemma}

\begin{proof}
For every integer $q\geq0$ there is a constant $C_q$ such that
\[
 \int_A^\infty v^q\e^{-v^2/2}\,\dd v
 \leq C_q(1+A^q)\e^{-A^2/4}
 \qquad(A\geq1).
\]
At $A=A_N(b)$,
\[
 \e^{-A_N(b)^2/4}
 =\e^{-8(N+1)\log b}
 =b^{-8(N+1)}.
\]
The polynomial factor in $A_N(b)$ is a power of $\log b$ and is absorbed
by the enormous surplus of inverse powers.  Sum over the finitely many
monomials of $P$.
\end{proof}

Combining \cref{prop:finite-exp-substitution,lem:gaussian-polynomial-tail}
with \eqref{eq:reference-contraction}, the central contour contributes
\begin{equation}
 \sum_{j=0}^{N}a_j(b)b^{-j}+\bigO(b^{-N-1}).
 \label{eq:central-mass-expansion}
\end{equation}

\subsection{The complementary contour}

The exact contour outside the central window is split into an intermediate
region and a far region.  The two estimates use different mechanisms.

\begin{proposition}[All-order minor-arc bounds; \Formalized]
\label{prop:minor-arcs}
For every fixed $N$, the normalized contour outside
$|v|\leq A_N(b)$ contributes $\bigO(b^{-N-1})$.
More precisely:
\begin{enumerate}
\item the intermediate region has a bound with fourfold power slack,
of order at most a constant times $b^{-4(N+1)}$ once the relevant
dyadic index is at least $b/4$;
\item the far region is bounded by a constant times
\[
 b\exp\!\left(-\frac{\log2}{4}b\right),
\]
which is flat against every inverse power of $b$.
\end{enumerate}
\end{proposition}

\begin{proof}
On the intermediate region the modulus estimate for the saddle kernel
contains a decay factor whose exponent is quadratic in the enlarged
radius.  Substitution of \eqref{eq:central-radius} yields
\[
 \exp\{-4(N+1)\log b\}=b^{-4(N+1)}
\]
after the standard comparison between the dyadic index and $b$.  The
remaining geometric and polynomial factors consume far less than the
three spare blocks of inverse powers.

In the far region the dyadic product supplies geometric decay in the
index.  Summing the geometric tail gives a bound of the displayed
exponential form.  For every $K$,
\[
 b^K\,b\e^{-(\log2)b/4}\longrightarrow0,
\]
so this contribution is $\bigO(b^{-N-1})$ for the requested fixed $N$.
\end{proof}

\subsection{The all-orders mass expansion}

\begin{theorem}[Normalized mass expansion; \Formalized]
\label{thm:mass-all-orders}
For every fixed $N\geq0$, as $t\to+\infty$ with
$b=\lambda(2^{-t})$,
\begin{equation}
 \mathcal G(2^{-t})
 =\sum_{j=0}^{N}\frac{a_j(b)}{b^j}
  +\bigO(b^{-N-1}). \label{eq:mass-all-orders}
\end{equation}
The coefficient functions are evaluated at the exact phase $b$.
\end{theorem}

\begin{proof}
The central part is \eqref{eq:central-mass-expansion}.  Add the
intermediate and far parts controlled by \cref{prop:minor-arcs}.  All
three errors are $\bigO(b^{-N-1})$.
\end{proof}

\subsection{Transfer through the logarithm}

\begin{theorem}[Logarithmic expansion on the dyadic real scale;
\Formalized]
\label{thm:log-dyadic-all-orders}
For every fixed $N\geq0$,
\begin{equation}
 \log\mathcal G(2^{-t})
 =\sum_{j=1}^{N}\frac{A_j(b)}{b^j}
  +\bigO(b^{-N-1}),
 \qquad b=\lambda(2^{-t}). \label{eq:log-mass-all-orders}
\end{equation}
\end{theorem}

\begin{proof}
The boundedness of the $a_j$ and \eqref{eq:mass-all-orders} with $N=0$
give $\mathcal G(2^{-t})=1+\bigO(b^{-1})$, hence
$\mathcal G>0$ and $|\mathcal G-1|<1/2$ eventually.  Taylor's theorem for
$\log(1+z)$ is therefore uniform.  Substitute
\eqref{eq:mass-all-orders} into the finite Taylor polynomial for the
logarithm.  The coefficient of $b^{-j}$ is exactly $A_j$, either by
\eqref{eq:A-compositions} or by the recursion
\eqref{eq:A-recursion}.  The logarithmic Taylor remainder and the
original mass remainder are both $\bigO(b^{-N-1})$.
\end{proof}

\subsection{Adding the flat tail and returning to \texorpdfstring{\(x\downarrow0\)}{x to 0}}

The tail $\mathcal T$ in \eqref{eq:exact-saddle-reduction} satisfies, for
every $K$,
\[
 \mathcal T(x)=\bigO(\lambda(x)^{-K}).
\]
It therefore does not change any coefficient in a Poincar\'e expansion.

\begin{theorem}[Full exact-phase expansion; \Formalized]
\label{thm:full-all-orders}
For every $N\geq0$,
\begin{equation}
 \log F(x)
 =\Main(x)+
   \sum_{j=1}^{N}\frac{A_j(\lambda(x))}{\lambda(x)^j}
   +\bigO(\lambda(x)^{-N-1})
 \qquad(x\downarrow0). \label{eq:full-all-orders}
\end{equation}
Equivalently, using a truncation with indices $j<N$,
\[
 \log F(x)-\Main(x)-
 \sum_{j=1}^{N-1}\frac{A_j(\lambda(x))}{\lambda(x)^j}
 =\bigO(\lambda(x)^{-N}).
\]
\end{theorem}

\begin{proof}
First combine \eqref{eq:exact-saddle-reduction},
\eqref{eq:log-mass-all-orders}, and the flatness of $\mathcal T$ on the
scale $x=2^{-t}$.  Next use the exact change of variables
$t=-\log_2x$.  As $x\downarrow0$, $t\to+\infty$ and
$\lambda(2^{-t})=\lambda(x)$, so the dyadic-real estimate transports
without approximation.  This yields \eqref{eq:full-all-orders}.
\end{proof}

\begin{warning}[Poincar\'e does not mean convergent]
The theorem is fixed-order.  For each $N$ there is an error constant
depending on $N$, and the central radius itself depends on $N$.  No
uniform control as $N\to\infty$ is asserted, and no convergence of the
infinite coefficient series is implied.
\end{warning}

\section{A proof-oriented concordance with the Lean modules}
\label{sec:module-concordance}

The following table records where the two missing clusters are proved.
Paths are relative to
\begin{center}
\path{Analysis/FabiusFunction/Lean/FabiusFunction/}.
\end{center}

\begingroup
\small
\begin{longtable}{>{\raggedright\arraybackslash}p{0.31\textwidth}
                  >{\raggedright\arraybackslash}p{0.49\textwidth}
                  >{\raggedright\arraybackslash}p{0.10\textwidth}}
\caption{Module concordance for this companion.}
\label{tab:module-concordance}\\
\toprule
Module & Principal role & Status\\
\midrule
\endfirsthead
\toprule
Module & Principal role & Status\\
\midrule
\endhead
\texttt{FabiusInverse.lean} &
Order isomorphism of $[0,1]$, totalized inverse, global continuity and
monotonicity, comparison equivalences, tails, reflection, exact dyadic
inversion, reciprocal first and second derivatives, interior smoothness,
strict closed-half curvature, transported flatness, effective steepness, and
endpoint quotient divergence. & \Formalized\\
\texttt{FabiusSaddleJet}\\\texttt{ClosedForm.lean} &
Operator solution of the recursive periodic jets and the harmonic-number
constant term. & \Formalized\\
\texttt{FabiusSaddleJet}\\\texttt{Stirling.lean} &
Identification of the operator coefficients with
$s(n+1,m+1)$ and the corresponding closed derivative formula. &
\Formalized\\
\texttt{FabiusSaddleExponent}\\\texttt{ClosedForm.lean} &
Collapse of each centered exponent coefficient to one bounded jet
monomial plus one universal monomial; parity. & \Formalized\\
\texttt{FabiusSaddleExpansion}\\\texttt{Coefficients.lean} &
Formal exponential coefficients, Gaussian contraction, mass
coefficients, logarithmic transfer, periodicity, smoothness, and
boundedness. & \Formalized\\
\texttt{FabiusSaddleLeading}\\\texttt{Coefficient.lean} &
Affine dependence of $a_j$ on the newest jet and coefficient
$(-1)^j/(2^jj!)$. & \Formalized\\
\texttt{FabiusSecondSaddle}\\\texttt{Correction.lean} &
Order-four exponential identity, explicit $B_4$, Gaussian contraction,
$a_2$, $A_2$, and regularity of the second correction. & \Formalized\\
\texttt{FabiusSecondSaddle}\\\texttt{Expansion.lean} &
Insertion of $A_2$ into the exact-phase asymptotic with a
third-order remainder. & \Formalized\\
\texttt{FabiusSaddleCentral}\\\texttt{AllOrders.lean} &
Exact central decomposition, finite reference polynomial, odd-term
cancellation, and integrated central remainders. & \Formalized\\
\texttt{FabiusSaddleTail}\\\texttt{AllOrders.lean} &
Order-dependent radius and complementary intermediate/far contour bounds
of arbitrary fixed inverse-power order. & \Formalized\\
\texttt{FabiusSaddleMass}\\\texttt{AllOrders.lean} &
Assembly of central Taylor control, finite exponential substitution,
whole-line Gaussian contraction, and minor arcs into the complete mass
expansion. & \Formalized\\
\texttt{FabiusFullAsymptotic}\\\texttt{Expansion.lean} &
Transfer from mass to logarithm, addition of the flat forward tail, and
transport from $x=2^{-t}$ to the full limit $x\downarrow0$. & \Formalized\\
\bottomrule
\end{longtable}
\endgroup

\subsection{Named declarations worth knowing}

The inverse interface is centered around
\begin{lstlisting}[style=pseudo]
fabiusOrderIso
fabiusReal_fabiusInv
fabiusInv_fabiusReal
strictMonoOn_fabiusInv
monotone_fabiusInv
continuous_fabiusInv
fabiusInv_le_iff_le_fabiusReal
fabiusReal_le_iff_le_fabiusInv
fabiusInv_lt_iff_lt_fabiusReal
fabiusReal_lt_iff_lt_fabiusInv
fabiusInv_one_sub
fabiusInv_fabiusDyadicUnit
le_two_pow_mul_fabiusInv_pow
le_two_pow_div_factorial_mul_fabiusInv_pow
mul_lt_fabiusInv
tendsto_fabiusInv_div_atTop
tendsto_one_sub_fabiusInv_div_one_sub_atTop
\end{lstlisting}
The all-orders endpoint interface culminates in
\begin{lstlisting}[style=pseudo]
fabiusSaddleKernelMass_dyadicLambert_hasAsymptoticExpansion
fabiusSaddleRatio_dyadicLambert_hasAsymptoticExpansion
log_fabius_dyadicReal_sub_sharpLambertMain_hasAsymptoticExpansion
log_fabius_sub_sharpLambertMain_hasAsymptoticExpansion
log_fabius_sub_sharpLambertExpansion_isBigO
\end{lstlisting}

These names are not needed to read the mathematics, but they make the
article useful as a map back into the source tree.

\section{What is formalized and what is newly synthesized}
\label{sec:status-ledger}

For clarity, the principal statements can be divided into three groups.

\subsection{Directly formalized}

The following content is represented directly by declarations in the
repository:
\begin{itemize}
\item the global totalized inverse package and all four comparison
equivalences;
\item the constant tails, reflection, midpoint, and exact dyadic
inversion;
\item interior $C^\infty$ regularity, the exact formulas for $\InvF'$ and
$\InvF''$, and strict concavity on $[0,1/2]$ and strict convexity on
$[1/2,1]$;
\item the exact analytic locus of the totalized inverse;
\item the effective and sharp inverse power inequalities, effective
linear steepness, and the two endpoint quotient limits;
\item the closed jet formulas, shifted Stirling conversion, and closed
exponent coefficients;
\item the formal exponential and logarithmic coefficient machinery,
top-jet coefficient, first and second corrections;
\item the arbitrary-order central, tail, mass, and logarithmic
asymptotic theorems.
\end{itemize}

\subsection{Classical consequences derived in this article}

The following statements are proved here from the formalized input:
\begin{itemize}
\item the exact smooth locus of the totalized inverse;
\item divergence of $\InvF(y)/y^\alpha$ for every $\alpha>0$;
\item the sharp equivalent \eqref{eq:sharp-inverse} and the scale
hierarchy \eqref{eq:power-below-G}--\eqref{eq:G-below-log};
\item the explicit all-$j$ composition sum \eqref{eq:mass-composition}
and the highest-$\Psi$-derivative corollary
\eqref{eq:highest-derivative}.
\end{itemize}

\subsection{Claims deliberately not made}

This companion does \emph{not} claim:
\begin{itemize}
\item convergence or Borel summability of the all-orders series;
\item an all-orders substitution of an approximate Lambert phase inside
the periodic coefficients;
\item a closed form for the inverse at arbitrary rational arguments;
\item analyticity of either $F$ or $\InvF$ inside the transition interval;
\item uniformity of the remainder constants as the truncation order tends
to infinity.
\end{itemize}

\section{Concluding synthesis}

The inverse function and the saddle coefficient machinery illuminate
opposite sides of the same endpoint phenomenon.

On the geometric side, $F$ approaches zero more rapidly than every
power.  Its inverse must therefore leave zero more rapidly than every
power, and the formalized comparison calculus turns this slogan into
effective inequalities with exact constants.  The sharp Lambert
expansion goes further: the inverse lives on the
stretched-logarithmic scale
\[
 \InvF(y)\asymp
 \sqrt{\log(1/y)}
 \exp\!\left(-\sqrt{2\log2\,\log(1/y)}\right),
\]
with the exact leading constant $1/(\e\sqrt{\log2})$.  Reflection transfers
the entire picture to the endpoint $1$.

On the asymptotic side, the apparent complexity of the periodic
corrections is organized by three simple structures.  First, the
recursive jets are solved by a product of commuting first-order
operators; the forcing terms sum to a harmonic number and the derivative
coefficients are shifted signed Stirling numbers.  Second, centering
collapses every exponent coefficient to two monomials, forcing parity and
eliminating all odd half-powers under Gaussian contraction.  Third, the
newest jet can enter only linearly, so the coefficient of the highest
derivative $\Psi^{(2j)}$ is known at every order before the lower-order
algebra is performed.

The analytic remainder proof then explains why these formal coefficients
are not merely symbolic.  The growing window
$A_N(b)=\sqrt{32(N+1)\log b}$ separates the local and global tasks:
locally it is still narrow enough for uniform Taylor expansion, while
globally it is wide enough to make both Gaussian and contour tails smaller
than the requested inverse power.  Finite exponential substitution and
logarithmic transfer complete a genuine Poincar\'e expansion at the exact
Lambert phase.

Together, these missing parts close two expository loops.  Flatness of
$F$ is now matched by a precise theory of the steep inverse, and the
headline all-orders expansion is now matched by a transparent account of
the coefficient algebra and the analytic estimates that make it true.

\appendix

\section{A compact inverse reference}
\label{app:inverse-reference}

For convenience, here is the inverse theory in one place.  For
$\InvF(y)=F_I^{-1}(\clamp(y))$:
\begin{align*}
&0\leq\InvF(y)\leq1,\qquad
 \InvF(y)=0\ (y\leq0),\qquad
 \InvF(y)=1\ (y\geq1),\\
&F(\InvF(y))=y\quad(0\leq y\leq1),\qquad
 \InvF(F(x))=x\quad(0\leq x\leq1),\\
&\InvF(1-y)=1-\InvF(y),\qquad
 \InvF(1/2)=1/2,\\
&\InvF(y)\leq x\iff y\leq F(x),\qquad
 F(x)\leq y\iff x\leq\InvF(y),\\
&\InvF(y)<x\iff y<F(x),\qquad
 F(x)<y\iff x<\InvF(y),\\
&\InvF(\mathsf F_n(a))
 =\min\{a,2^n\}/2^n,\\
&y\leq C_n\InvF(y)^n,\qquad
 y\leq C_n\InvF(y)^n/n!
 \quad\text{when }2^n\InvF(y)\leq1,\\
&\InvF(y)/y\to+\infty\quad(y\downarrow0),\qquad
 (1-\InvF(y))/(1-y)\to+\infty\quad(y\uparrow1),\\
&\InvF(y)\sim
 \frac{\sqrt{\log(1/y)}}{\e\sqrt{\log2}}\,
 \exp\!\left[-\sqrt{2\log2\,\log(1/y)}\right].
\end{align*}

\section{A compact coefficient reference}
\label{app:coefficient-reference}

With $L=\log2$,
\begin{align*}
p_0&=\frac12+\frac{\Psi'}L,\\
p_{n+1}&=(D/L-(n+1))p_n+\frac{(-1)^{n+1}n!}{L},\\
d_n&=p_n+(-1)^{n+1}n!,\\
c(n,m)&=s(n+1,m+1),\\
d_n&=(-1)^nn!\left(\frac{H_n}{L}-\frac12\right)
 +\sum_{m=0}^nc(n,m)\frac{\Psi^{(m+1)}}{L^{m+1}},\\
E_m&=\ii^m\left(\frac{d_{m-1}}{m!}v^m+
          \frac{(-1)^{m+1}}{m+2}v^{m+2}\right),\\
nB_n&=\sum_{m=1}^n mE_mB_{n-m},\\
a_j&=\Gauss[B_{2j}],\\
A_n&=a_n-\frac1n\sum_{k=1}^{n-1}kA_ka_{n-k},\\
[d_{2j-1}]\,A_j&=\frac{(-1)^j}{2^jj!},\\
[\Psi^{(2j)}]\,A_j&=\frac{(-1)^j}{2^jj!L^{2j}}.
\end{align*}
The first two logarithmic coefficients are
\begin{align*}
A_1={}&-\frac{d_0^2}{2}-d_0-\frac{d_1}{2}-\frac1{12},\\
A_2={}&\frac1{24}\bigl(
8d_0^3+12d_0^2d_1+12d_0^2+48d_0d_1+12d_0d_2\\
&\hspace{4.1em}{}+6d_1^2+24d_1+20d_2+3d_3
\bigr).
\end{align*}

\section{Repository paths used in the audit}
\label{app:paths}

\begin{lstlisting}[style=pseudo]
Analysis/FabiusFunction/
  Lean/FabiusFunction/FabiusInverse.lean
  Lean/FabiusFunction/FabiusSaddleJetClosedForm.lean
  Lean/FabiusFunction/FabiusSaddleJetStirling.lean
  Lean/FabiusFunction/FabiusSaddleExponentClosedForm.lean
  Lean/FabiusFunction/FabiusSaddleLeadingCoefficient.lean
  Lean/FabiusFunction/FabiusSecondSaddleCorrection.lean
  Lean/FabiusFunction/FabiusSecondSaddleExpansion.lean
  Lean/FabiusFunction/FabiusSaddleCentralAllOrders.lean
  Lean/FabiusFunction/FabiusSaddleTailAllOrders.lean
  Lean/FabiusFunction/FabiusSaddleMassAllOrders.lean
  Lean/FabiusFunction/FabiusFullAsymptoticExpansion.lean
  docs/Fabius_Function_and_Rvachev_Up/
    Fabius_Function_and_Rvachev_Up.tex
\end{lstlisting}

\begin{thebibliography}{99}

\bibitem{fabius1966}
J.~Fabius,
\newblock A probabilistic example of a nowhere analytic $C^\infty$-function,
\newblock \emph{Zeitschrift f\"ur Wahrscheinlichkeitstheorie und Verwandte
Gebiete} \textbf{5} (1966), 173--174.
\newblock DOI: \nolinkurl{10.1007/BF00536652}.

\bibitem{rvachev1990}
V.~A.~Rvachev,
\newblock Compactly supported solutions of functional-differential equations
and their applications,
\newblock \emph{Russian Mathematical Surveys} \textbf{45} (1990), no.~1,
87--120.

\bibitem{arias1982}
J.~Arias de Reyna,
\newblock Definici\'on y estudio de una funci\'on indefinidamente
diferenciable de soporte compacto,
\newblock \emph{Revista de la Real Academia de Ciencias Exactas, F\'isicas y
Naturales de Madrid} \textbf{76} (1982), no.~1, 21--38.
\newblock English translation:
\emph{An infinitely differentiable function with compact support:
Definition and properties}, arXiv:\nolinkurl{1702.05442} (2017).

\bibitem{arias2018}
J.~Arias de Reyna,
\newblock Arithmetic of the Fabius function,
\newblock \emph{INTEGERS} \textbf{18} (2018), Article A51;
arXiv:\nolinkurl{1702.06487}, version~3.

\bibitem{corless1996}
R.~M.~Corless, G.~H.~Gonnet, D.~E.~G.~Hare, D.~J.~Jeffrey, and D.~E.~Knuth,
\newblock On the Lambert $W$ function,
\newblock \emph{Advances in Computational Mathematics}
\textbf{5} (1996), 329--359.

\bibitem{proveit}
V.~Reshetnikov,
\newblock \emph{ProveIt: Analysis/FabiusFunction},
\newblock Lean~4 formal development, repository snapshot 25 August 2026.
\newblock
\url{https://github.com/VladimirReshetnikov/ProveIt/tree/main/Analysis/FabiusFunction}.

\bibitem{maincompanion}
V.~Reshetnikov,
\newblock \emph{The Fabius Function and Rvachev's Up-Function},
\newblock repository document in
\path{Analysis/FabiusFunction/docs/Fabius_Function_and_Rvachev_Up},
snapshot 25 August 2026.

\bibitem{smallargument}
V.~Reshetnikov,
\newblock \emph{The Small-Argument Asymptotic Expansion of the Fabius
Function},
\newblock technical repository document accompanying the Lean saddle
development, snapshot 25 August 2026.

\end{thebibliography}

\end{document}
